% ============================================================
% Flujo Operativo del Ciclo Completo de Matrícula — Cap. 4.1.3
% 5 Fases correlativas (F1 Preinscripción → F5 Asignación)
% Proyecto BTH V2.0 — MTP San Julián
% Formato: TikZ flowchart vertical con decisión opcional Fase 2 — v2 espaciado ampliado
% Uso: % ============================================================
% Flujo Operativo del Ciclo Completo de Matrícula — Cap. 4.1.3
% 5 Fases correlativas (F1 Preinscripción → F5 Asignación)
% Proyecto BTH V2.0 — MTP San Julián
% Formato: TikZ flowchart vertical con decisión opcional Fase 2 — v2 espaciado ampliado
% Uso: % ============================================================
% Flujo Operativo del Ciclo Completo de Matrícula — Cap. 4.1.3
% 5 Fases correlativas (F1 Preinscripción → F5 Asignación)
% Proyecto BTH V2.0 — MTP San Julián
% Formato: TikZ flowchart vertical con decisión opcional Fase 2 — v2 espaciado ampliado
% Uso: % ============================================================
% Flujo Operativo del Ciclo Completo de Matrícula — Cap. 4.1.3
% 5 Fases correlativas (F1 Preinscripción → F5 Asignación)
% Proyecto BTH V2.0 — MTP San Julián
% Formato: TikZ flowchart vertical con decisión opcional Fase 2 — v2 espaciado ampliado
% Uso: \input{imagenes/flujo_matricula} dentro de figure
% ============================================================
\begin{tikzpicture}[
    fase/.style={
        rectangle, rounded corners=4pt,
        draw=primary!75, fill=white,
        text=darktext, align=center,
        font=\scriptsize,
        text width=9.8cm, minimum height=1.05cm,
        inner sep=4pt
    },
    fasePrimary/.style={
        rectangle, rounded corners=4pt,
        draw=primary, fill=primary,
        text=white, align=center,
        font=\scriptsize\bfseries,
        text width=9.8cm, minimum height=1.05cm,
        inner sep=4pt
    },
    faseLight/.style={
        rectangle, rounded corners=4pt,
        draw=primary!55, fill=secondary!12,
        text=darktext, align=center,
        font=\scriptsize,
        text width=9.8cm, minimum height=1.05cm,
        inner sep=4pt
    },
    faseBranch/.style={
        rectangle, rounded corners=4pt,
        draw=primary!75, fill=white,
        text=darktext, align=center,
        font=\scriptsize,
        text width=4.4cm, minimum height=1.15cm,
        inner sep=4pt
    },
    faseBranchLight/.style={
        rectangle, rounded corners=4pt,
        draw=primary!55, fill=secondary!12,
        text=darktext, align=center,
        font=\scriptsize,
        text width=4.4cm, minimum height=1.15cm,
        inner sep=4pt
    },
    decision/.style={
        diamond, aspect=2.4,
        draw=primary!80, fill=lightgray,
        text=darktext, align=center,
        font=\scriptsize\bfseries,
        inner sep=3pt, text width=2.6cm
    },
    terminator/.style={
        ellipse,
        draw=primary!75, fill=primary!12,
        text=darktext, align=center,
        font=\scriptsize\bfseries,
        minimum width=2.8cm, minimum height=0.75cm
    },
    flecha/.style={
        -{Stealth[length=2.8mm,width=2.1mm]},
        thick, draw=primary!85,
        rounded corners=2pt
    },
    flechaDashed/.style={
        -{Stealth[length=2.6mm,width=2mm]},
        thick, draw=primary!55, dashed,
        rounded corners=2pt
    },
    etiqueta/.style={
        font=\tiny, text=primary!80, fill=white, inner sep=1pt
    },
    nota/.style={
        font=\tiny\itshape, text=darktext!65, align=center
    },
    actor/.style={
        font=\tiny\bfseries, text=primary!65, align=center
    }
]

% --- Actores superiores (sutil) — espaciado ampliado v4 ---
\node[actor] at (-5.70, 6.70) {POSTULANTE};
\node[actor] at (0, 6.70) {SISTEMA MVC};
\node[actor] at (5.70, 6.70) {SECRETARÍA BTH};
\draw[primary!18, thick] (-6.95, 6.48) -- (6.95, 6.48);

% --- INICIO — gap +0.40cm ---
\node[terminator] (inicio) at (0, 6.00) {INICIO};

% --- PASO 1 ---
\node[fasePrimary] (f1) at (0, 4.75) {\textbf{Paso 1. Preinscripción y reserva}\\ \footnotesize\normalfont Formulario revisado \textbullet{} Elige 7 carreras \textbullet{} Guarda cupo seguro};

% Flecha inicio -> F1 (más aire)
\draw[flecha] (inicio.south) -- (f1.north);
\node[etiqueta, rotate=90] at (-6.05, 5.38) {móvil/PC};

% --- DECISIÓN Paso 2 ---
\node[decision] (dec2) at (0, 3.10) {¿Olvidó\\contraseña?};

\draw[flecha] (f1.south) -- (dec2.north);

% Rama SI — Recuperación
\node[faseBranch] (f2si) at (-5.70, 1.35) {\textbf{Paso 2. Recuperar}\\ \footnotesize\normalfont Con carnet y fecha};

% Rama NO
\node[faseBranchLight] (f2no) at (5.70, 1.35) {\textbf{Paso 2 (No): Sigue}\\ \footnotesize\normalfont Sin recuperar};

\draw[flecha] (dec2.west) -- ++(-0.90,0) |- (f2si.north);
\node[etiqueta] at (-2.90, 2.35) {Sí};
\draw[flecha] (dec2.east) -- ++(0.90,0) |- (f2no.north);
\node[etiqueta] at (2.90, 2.35) {No};

% Convergencia hacia Fase 3 — gap ampliado +0.70cm, codos suavizados
\coordinate (conv) at (0, -0.35);
\draw[flechaDashed] (f2si.south) -- (-5.70, -0.35) -- (conv);
\draw[flechaDashed] (f2no.south) -- (5.70, -0.35) -- (conv);

% --- PASO 3 ---
\node[faseLight] (f3) at (0, -1.45) {\textbf{Paso 3. Comprobante}\\ \footnotesize\normalfont PDF con código único};

\draw[flecha] (conv) -- (f3.north);

% --- PASO 4 ---
\node[fase] (f4) at (0, -3.20) {\textbf{Paso 4. Entrega de documentos}\\ \footnotesize\normalfont Revisa libretas, carnet y certificado};

\draw[flecha] (f3.south) -- (f4.north);
\node[etiqueta, rotate=90] at (-6.05, -2.33) {gestión siguiente};

% --- PASO 5 ---
\node[fasePrimary] (f5) at (0, -4.90) {\textbf{Paso 5. Ubicación en cursos}\\ \footnotesize\normalfont Por colegio → 4to A / 4to B, listas finales};

\draw[flecha] (f4.south) -- (f5.north);

% --- FIN — gap +0.75cm ---
\node[terminator] (fin) at (0, -6.40) {FIN};

\draw[flecha] (f5.south) -- (fin.north);

% --- Línea lateral actor (vertical sutil, extendida v4) ---
\draw[primary!15, line width=1.8pt] (-6.85, 6.10) -- (-6.85, -6.50);
\draw[primary!15, line width=1.8pt] (6.85, 6.10) -- (6.85, -6.50);

% --- Notas inferiores (sencillo) ---
\node[nota] at (0, -7.05) {Ciclo completo: final de una gestión (Pasos 1--3) $\rightarrow$ inicio de la siguiente (Pasos 4--5)};
\node[nota, text=darktext!50] at (0, -7.40) {Tiempo: 3--5 min por internet vs. 1--3 días de fila \textbullet{} Atención 24 h \textbullet{} Todo queda registrado};

\end{tikzpicture}
 dentro de figure
% ============================================================
\begin{tikzpicture}[
    fase/.style={
        rectangle, rounded corners=4pt,
        draw=primary!75, fill=white,
        text=darktext, align=center,
        font=\scriptsize,
        text width=9.8cm, minimum height=1.05cm,
        inner sep=4pt
    },
    fasePrimary/.style={
        rectangle, rounded corners=4pt,
        draw=primary, fill=primary,
        text=white, align=center,
        font=\scriptsize\bfseries,
        text width=9.8cm, minimum height=1.05cm,
        inner sep=4pt
    },
    faseLight/.style={
        rectangle, rounded corners=4pt,
        draw=primary!55, fill=secondary!12,
        text=darktext, align=center,
        font=\scriptsize,
        text width=9.8cm, minimum height=1.05cm,
        inner sep=4pt
    },
    faseBranch/.style={
        rectangle, rounded corners=4pt,
        draw=primary!75, fill=white,
        text=darktext, align=center,
        font=\scriptsize,
        text width=4.4cm, minimum height=1.15cm,
        inner sep=4pt
    },
    faseBranchLight/.style={
        rectangle, rounded corners=4pt,
        draw=primary!55, fill=secondary!12,
        text=darktext, align=center,
        font=\scriptsize,
        text width=4.4cm, minimum height=1.15cm,
        inner sep=4pt
    },
    decision/.style={
        diamond, aspect=2.4,
        draw=primary!80, fill=lightgray,
        text=darktext, align=center,
        font=\scriptsize\bfseries,
        inner sep=3pt, text width=2.6cm
    },
    terminator/.style={
        ellipse,
        draw=primary!75, fill=primary!12,
        text=darktext, align=center,
        font=\scriptsize\bfseries,
        minimum width=2.8cm, minimum height=0.75cm
    },
    flecha/.style={
        -{Stealth[length=2.8mm,width=2.1mm]},
        thick, draw=primary!85,
        rounded corners=2pt
    },
    flechaDashed/.style={
        -{Stealth[length=2.6mm,width=2mm]},
        thick, draw=primary!55, dashed,
        rounded corners=2pt
    },
    etiqueta/.style={
        font=\tiny, text=primary!80, fill=white, inner sep=1pt
    },
    nota/.style={
        font=\tiny\itshape, text=darktext!65, align=center
    },
    actor/.style={
        font=\tiny\bfseries, text=primary!65, align=center
    }
]

% --- Actores superiores (sutil) — espaciado ampliado v4 ---
\node[actor] at (-5.70, 6.70) {POSTULANTE};
\node[actor] at (0, 6.70) {SISTEMA MVC};
\node[actor] at (5.70, 6.70) {SECRETARÍA BTH};
\draw[primary!18, thick] (-6.95, 6.48) -- (6.95, 6.48);

% --- INICIO — gap +0.40cm ---
\node[terminator] (inicio) at (0, 6.00) {INICIO};

% --- PASO 1 ---
\node[fasePrimary] (f1) at (0, 4.75) {\textbf{Paso 1. Preinscripción y reserva}\\ \footnotesize\normalfont Formulario revisado \textbullet{} Elige 7 carreras \textbullet{} Guarda cupo seguro};

% Flecha inicio -> F1 (más aire)
\draw[flecha] (inicio.south) -- (f1.north);
\node[etiqueta, rotate=90] at (-6.05, 5.38) {móvil/PC};

% --- DECISIÓN Paso 2 ---
\node[decision] (dec2) at (0, 3.10) {¿Olvidó\\contraseña?};

\draw[flecha] (f1.south) -- (dec2.north);

% Rama SI — Recuperación
\node[faseBranch] (f2si) at (-5.70, 1.35) {\textbf{Paso 2. Recuperar}\\ \footnotesize\normalfont Con carnet y fecha};

% Rama NO
\node[faseBranchLight] (f2no) at (5.70, 1.35) {\textbf{Paso 2 (No): Sigue}\\ \footnotesize\normalfont Sin recuperar};

\draw[flecha] (dec2.west) -- ++(-0.90,0) |- (f2si.north);
\node[etiqueta] at (-2.90, 2.35) {Sí};
\draw[flecha] (dec2.east) -- ++(0.90,0) |- (f2no.north);
\node[etiqueta] at (2.90, 2.35) {No};

% Convergencia hacia Fase 3 — gap ampliado +0.70cm, codos suavizados
\coordinate (conv) at (0, -0.35);
\draw[flechaDashed] (f2si.south) -- (-5.70, -0.35) -- (conv);
\draw[flechaDashed] (f2no.south) -- (5.70, -0.35) -- (conv);

% --- PASO 3 ---
\node[faseLight] (f3) at (0, -1.45) {\textbf{Paso 3. Comprobante}\\ \footnotesize\normalfont PDF con código único};

\draw[flecha] (conv) -- (f3.north);

% --- PASO 4 ---
\node[fase] (f4) at (0, -3.20) {\textbf{Paso 4. Entrega de documentos}\\ \footnotesize\normalfont Revisa libretas, carnet y certificado};

\draw[flecha] (f3.south) -- (f4.north);
\node[etiqueta, rotate=90] at (-6.05, -2.33) {gestión siguiente};

% --- PASO 5 ---
\node[fasePrimary] (f5) at (0, -4.90) {\textbf{Paso 5. Ubicación en cursos}\\ \footnotesize\normalfont Por colegio → 4to A / 4to B, listas finales};

\draw[flecha] (f4.south) -- (f5.north);

% --- FIN — gap +0.75cm ---
\node[terminator] (fin) at (0, -6.40) {FIN};

\draw[flecha] (f5.south) -- (fin.north);

% --- Línea lateral actor (vertical sutil, extendida v4) ---
\draw[primary!15, line width=1.8pt] (-6.85, 6.10) -- (-6.85, -6.50);
\draw[primary!15, line width=1.8pt] (6.85, 6.10) -- (6.85, -6.50);

% --- Notas inferiores (sencillo) ---
\node[nota] at (0, -7.05) {Ciclo completo: final de una gestión (Pasos 1--3) $\rightarrow$ inicio de la siguiente (Pasos 4--5)};
\node[nota, text=darktext!50] at (0, -7.40) {Tiempo: 3--5 min por internet vs. 1--3 días de fila \textbullet{} Atención 24 h \textbullet{} Todo queda registrado};

\end{tikzpicture}
 dentro de figure
% ============================================================
\begin{tikzpicture}[
    fase/.style={
        rectangle, rounded corners=4pt,
        draw=primary!75, fill=white,
        text=darktext, align=center,
        font=\scriptsize,
        text width=9.8cm, minimum height=1.05cm,
        inner sep=4pt
    },
    fasePrimary/.style={
        rectangle, rounded corners=4pt,
        draw=primary, fill=primary,
        text=white, align=center,
        font=\scriptsize\bfseries,
        text width=9.8cm, minimum height=1.05cm,
        inner sep=4pt
    },
    faseLight/.style={
        rectangle, rounded corners=4pt,
        draw=primary!55, fill=secondary!12,
        text=darktext, align=center,
        font=\scriptsize,
        text width=9.8cm, minimum height=1.05cm,
        inner sep=4pt
    },
    faseBranch/.style={
        rectangle, rounded corners=4pt,
        draw=primary!75, fill=white,
        text=darktext, align=center,
        font=\scriptsize,
        text width=4.4cm, minimum height=1.15cm,
        inner sep=4pt
    },
    faseBranchLight/.style={
        rectangle, rounded corners=4pt,
        draw=primary!55, fill=secondary!12,
        text=darktext, align=center,
        font=\scriptsize,
        text width=4.4cm, minimum height=1.15cm,
        inner sep=4pt
    },
    decision/.style={
        diamond, aspect=2.4,
        draw=primary!80, fill=lightgray,
        text=darktext, align=center,
        font=\scriptsize\bfseries,
        inner sep=3pt, text width=2.6cm
    },
    terminator/.style={
        ellipse,
        draw=primary!75, fill=primary!12,
        text=darktext, align=center,
        font=\scriptsize\bfseries,
        minimum width=2.8cm, minimum height=0.75cm
    },
    flecha/.style={
        -{Stealth[length=2.8mm,width=2.1mm]},
        thick, draw=primary!85,
        rounded corners=2pt
    },
    flechaDashed/.style={
        -{Stealth[length=2.6mm,width=2mm]},
        thick, draw=primary!55, dashed,
        rounded corners=2pt
    },
    etiqueta/.style={
        font=\tiny, text=primary!80, fill=white, inner sep=1pt
    },
    nota/.style={
        font=\tiny\itshape, text=darktext!65, align=center
    },
    actor/.style={
        font=\tiny\bfseries, text=primary!65, align=center
    }
]

% --- Actores superiores (sutil) — espaciado ampliado v4 ---
\node[actor] at (-5.70, 6.70) {POSTULANTE};
\node[actor] at (0, 6.70) {SISTEMA MVC};
\node[actor] at (5.70, 6.70) {SECRETARÍA BTH};
\draw[primary!18, thick] (-6.95, 6.48) -- (6.95, 6.48);

% --- INICIO — gap +0.40cm ---
\node[terminator] (inicio) at (0, 6.00) {INICIO};

% --- PASO 1 ---
\node[fasePrimary] (f1) at (0, 4.75) {\textbf{Paso 1. Preinscripción y reserva}\\ \footnotesize\normalfont Formulario revisado \textbullet{} Elige 7 carreras \textbullet{} Guarda cupo seguro};

% Flecha inicio -> F1 (más aire)
\draw[flecha] (inicio.south) -- (f1.north);
\node[etiqueta, rotate=90] at (-6.05, 5.38) {móvil/PC};

% --- DECISIÓN Paso 2 ---
\node[decision] (dec2) at (0, 3.10) {¿Olvidó\\contraseña?};

\draw[flecha] (f1.south) -- (dec2.north);

% Rama SI — Recuperación
\node[faseBranch] (f2si) at (-5.70, 1.35) {\textbf{Paso 2. Recuperar}\\ \footnotesize\normalfont Con carnet y fecha};

% Rama NO
\node[faseBranchLight] (f2no) at (5.70, 1.35) {\textbf{Paso 2 (No): Sigue}\\ \footnotesize\normalfont Sin recuperar};

\draw[flecha] (dec2.west) -- ++(-0.90,0) |- (f2si.north);
\node[etiqueta] at (-2.90, 2.35) {Sí};
\draw[flecha] (dec2.east) -- ++(0.90,0) |- (f2no.north);
\node[etiqueta] at (2.90, 2.35) {No};

% Convergencia hacia Fase 3 — gap ampliado +0.70cm, codos suavizados
\coordinate (conv) at (0, -0.35);
\draw[flechaDashed] (f2si.south) -- (-5.70, -0.35) -- (conv);
\draw[flechaDashed] (f2no.south) -- (5.70, -0.35) -- (conv);

% --- PASO 3 ---
\node[faseLight] (f3) at (0, -1.45) {\textbf{Paso 3. Comprobante}\\ \footnotesize\normalfont PDF con código único};

\draw[flecha] (conv) -- (f3.north);

% --- PASO 4 ---
\node[fase] (f4) at (0, -3.20) {\textbf{Paso 4. Entrega de documentos}\\ \footnotesize\normalfont Revisa libretas, carnet y certificado};

\draw[flecha] (f3.south) -- (f4.north);
\node[etiqueta, rotate=90] at (-6.05, -2.33) {gestión siguiente};

% --- PASO 5 ---
\node[fasePrimary] (f5) at (0, -4.90) {\textbf{Paso 5. Ubicación en cursos}\\ \footnotesize\normalfont Por colegio → 4to A / 4to B, listas finales};

\draw[flecha] (f4.south) -- (f5.north);

% --- FIN — gap +0.75cm ---
\node[terminator] (fin) at (0, -6.40) {FIN};

\draw[flecha] (f5.south) -- (fin.north);

% --- Línea lateral actor (vertical sutil, extendida v4) ---
\draw[primary!15, line width=1.8pt] (-6.85, 6.10) -- (-6.85, -6.50);
\draw[primary!15, line width=1.8pt] (6.85, 6.10) -- (6.85, -6.50);

% --- Notas inferiores (sencillo) ---
\node[nota] at (0, -7.05) {Ciclo completo: final de una gestión (Pasos 1--3) $\rightarrow$ inicio de la siguiente (Pasos 4--5)};
\node[nota, text=darktext!50] at (0, -7.40) {Tiempo: 3--5 min por internet vs. 1--3 días de fila \textbullet{} Atención 24 h \textbullet{} Todo queda registrado};

\end{tikzpicture}
 dentro de figure
% ============================================================
\begin{tikzpicture}[
    fase/.style={
        rectangle, rounded corners=4pt,
        draw=primary!75, fill=white,
        text=darktext, align=center,
        font=\scriptsize,
        text width=9.8cm, minimum height=1.05cm,
        inner sep=4pt
    },
    fasePrimary/.style={
        rectangle, rounded corners=4pt,
        draw=primary, fill=primary,
        text=white, align=center,
        font=\scriptsize\bfseries,
        text width=9.8cm, minimum height=1.05cm,
        inner sep=4pt
    },
    faseLight/.style={
        rectangle, rounded corners=4pt,
        draw=primary!55, fill=secondary!12,
        text=darktext, align=center,
        font=\scriptsize,
        text width=9.8cm, minimum height=1.05cm,
        inner sep=4pt
    },
    faseBranch/.style={
        rectangle, rounded corners=4pt,
        draw=primary!75, fill=white,
        text=darktext, align=center,
        font=\scriptsize,
        text width=4.4cm, minimum height=1.15cm,
        inner sep=4pt
    },
    faseBranchLight/.style={
        rectangle, rounded corners=4pt,
        draw=primary!55, fill=secondary!12,
        text=darktext, align=center,
        font=\scriptsize,
        text width=4.4cm, minimum height=1.15cm,
        inner sep=4pt
    },
    decision/.style={
        diamond, aspect=2.4,
        draw=primary!80, fill=lightgray,
        text=darktext, align=center,
        font=\scriptsize\bfseries,
        inner sep=3pt, text width=2.6cm
    },
    terminator/.style={
        ellipse,
        draw=primary!75, fill=primary!12,
        text=darktext, align=center,
        font=\scriptsize\bfseries,
        minimum width=2.8cm, minimum height=0.75cm
    },
    flecha/.style={
        -{Stealth[length=2.8mm,width=2.1mm]},
        thick, draw=primary!85,
        rounded corners=2pt
    },
    flechaDashed/.style={
        -{Stealth[length=2.6mm,width=2mm]},
        thick, draw=primary!55, dashed,
        rounded corners=2pt
    },
    etiqueta/.style={
        font=\tiny, text=primary!80, fill=white, inner sep=1pt
    },
    nota/.style={
        font=\tiny\itshape, text=darktext!65, align=center
    },
    actor/.style={
        font=\tiny\bfseries, text=primary!65, align=center
    }
]

% --- Actores superiores (sutil) — espaciado ampliado v4 ---
\node[actor] at (-5.70, 6.70) {POSTULANTE};
\node[actor] at (0, 6.70) {SISTEMA MVC};
\node[actor] at (5.70, 6.70) {SECRETARÍA BTH};
\draw[primary!18, thick] (-6.95, 6.48) -- (6.95, 6.48);

% --- INICIO — gap +0.40cm ---
\node[terminator] (inicio) at (0, 6.00) {INICIO};

% --- PASO 1 ---
\node[fasePrimary] (f1) at (0, 4.75) {\textbf{Paso 1. Preinscripción y reserva}\\ \footnotesize\normalfont Formulario revisado \textbullet{} Elige 7 carreras \textbullet{} Guarda cupo seguro};

% Flecha inicio -> F1 (más aire)
\draw[flecha] (inicio.south) -- (f1.north);
\node[etiqueta, rotate=90] at (-6.05, 5.38) {móvil/PC};

% --- DECISIÓN Paso 2 ---
\node[decision] (dec2) at (0, 3.10) {¿Olvidó\\contraseña?};

\draw[flecha] (f1.south) -- (dec2.north);

% Rama SI — Recuperación
\node[faseBranch] (f2si) at (-5.70, 1.35) {\textbf{Paso 2. Recuperar}\\ \footnotesize\normalfont Con carnet y fecha};

% Rama NO
\node[faseBranchLight] (f2no) at (5.70, 1.35) {\textbf{Paso 2 (No): Sigue}\\ \footnotesize\normalfont Sin recuperar};

\draw[flecha] (dec2.west) -- ++(-0.90,0) |- (f2si.north);
\node[etiqueta] at (-2.90, 2.35) {Sí};
\draw[flecha] (dec2.east) -- ++(0.90,0) |- (f2no.north);
\node[etiqueta] at (2.90, 2.35) {No};

% Convergencia hacia Fase 3 — gap ampliado +0.70cm, codos suavizados
\coordinate (conv) at (0, -0.35);
\draw[flechaDashed] (f2si.south) -- (-5.70, -0.35) -- (conv);
\draw[flechaDashed] (f2no.south) -- (5.70, -0.35) -- (conv);

% --- PASO 3 ---
\node[faseLight] (f3) at (0, -1.45) {\textbf{Paso 3. Comprobante}\\ \footnotesize\normalfont PDF con código único};

\draw[flecha] (conv) -- (f3.north);

% --- PASO 4 ---
\node[fase] (f4) at (0, -3.20) {\textbf{Paso 4. Entrega de documentos}\\ \footnotesize\normalfont Revisa libretas, carnet y certificado};

\draw[flecha] (f3.south) -- (f4.north);
\node[etiqueta, rotate=90] at (-6.05, -2.33) {gestión siguiente};

% --- PASO 5 ---
\node[fasePrimary] (f5) at (0, -4.90) {\textbf{Paso 5. Ubicación en cursos}\\ \footnotesize\normalfont Por colegio → 4to A / 4to B, listas finales};

\draw[flecha] (f4.south) -- (f5.north);

% --- FIN — gap +0.75cm ---
\node[terminator] (fin) at (0, -6.40) {FIN};

\draw[flecha] (f5.south) -- (fin.north);

% --- Línea lateral actor (vertical sutil, extendida v4) ---
\draw[primary!15, line width=1.8pt] (-6.85, 6.10) -- (-6.85, -6.50);
\draw[primary!15, line width=1.8pt] (6.85, 6.10) -- (6.85, -6.50);

% --- Notas inferiores (sencillo) ---
\node[nota] at (0, -7.05) {Ciclo completo: final de una gestión (Pasos 1--3) $\rightarrow$ inicio de la siguiente (Pasos 4--5)};
\node[nota, text=darktext!50] at (0, -7.40) {Tiempo: 3--5 min por internet vs. 1--3 días de fila \textbullet{} Atención 24 h \textbullet{} Todo queda registrado};

\end{tikzpicture}
